% !Mode:: "TeX:UTF-8"
% !TEX program  = xelatex
\section{实验设计}
\subsection{实验目标}
本节旨在系统评估本文提出的基于预测的适配器调度方法在多 LoRA 推理场景下的有效性。我们主要关注以下几个问题：与传统缓存策略（如 LRU 和 LFU）相比，所提方法是否能够在有限显存约束下提升系统吞吐量并降低请求延迟；是否能够更高效地利用显存资源，从而提升热点适配器命中率；引入预测机制（EMA）对性能的具体贡献；以及在不同负载模式下的稳定性与泛化能力。

\subsection{实验设置}
使用 NVIDIA A800（80GB）显卡，系统为 Ubuntu Linux 5.4.0，CPU 为 Intel Xeon Gold 5320 @ 2.20GHz，内存为 692GB。模型方面，选用 Llama-2-13b 作为基础模型，基础模型占用显存约 28GB，并构建多个 LoRA 适配器模拟多任务推理场景。所有 LoRA 适配器总显存占用约 5GB，GPU 可用于适配器的显存预算设置为 2GB 或 4GB。

\subsection{工作负载}
\subsubsection{长尾分布}
为模拟真实在线服务中的长尾访问特征，采用 Zipf 分布生成请求。对于第 $i$ 个适配器 $a_i$，其访问概率定义为
\begin{equation}
P(x_t=a_i)=\frac{i^{-s}}{\sum_{j=1}^{M}j^{-s}},\quad s=1.0
\end{equation}
该工作负载用于评估策略在长尾热点分布下对高频适配器的保留能力。

\subsubsection{动态热点迁移}
请求序列被划分为多个阶段，每阶段包含近似相同数量请求。第 $p$ 阶段定义热点集合 $\mathcal{H}_p$，并设置热点请求比例 $\theta=0.8$。通过控制相邻阶段热点集合重叠比例 $\rho$，模拟热点迁移强度。该负载用于评估策略对非平稳访问模式的自适应能力。

\subsubsection{异构适配器}
根据适配器显存占用将其划分为小/中型集合 $S$ 与大型集合 $L$。设置 $P(x_t\in S)=\phi=0.75$，用于模拟高频小适配器与低频大适配器之间的显存竞争，评估策略的显存感知能力。

\subsection{实现细节}
系统由 AdapterRuntime 统一管理，基础模型常驻 GPU，LoRA 适配器由 AdapterPool 在显存预算内动态加载与淘汰。对比策略包括 FIFO、LRU、LFU 与本文方法，并支持消融实验。记录指标包括：命中率、未命中次数、淘汰次数、平均延迟、P95 延迟、吞吐量和总加载时间。

\subsection{数据分析}
\subsubsection{总体表现}
在 4GB 显存预算下，本文方法在长尾分布、热点偏移、异构分布三类负载中均达到更高命中率与更低平均服务延迟。在 2GB 和 4GB 两种预算对比下，本文方法均显著降低淘汰次数，且在紧显存（2GB）场景收益更明显。

\subsubsection{消融实验}
消融结果显示：去除显存归一化模块后命中率显著下降；去除短期频率项与成本感知项后性能变化相对较小。结果表明显存感知机制对提升系统整体资源利用效率至关重要。
